% swim-times-stp.tex
% Software Test Plan for the swim-times program
% ISO/IEC/IEEE 29119-3:2021 Test Plan structure
% Compiled with: pdflatex swim-times-stp.tex

\documentclass[12pt]{article}

\usepackage{geometry}
\geometry{letterpaper, margin=1in}

\usepackage{hyperref}
\usepackage{booktabs}
\usepackage{array}
\usepackage{longtable}
\usepackage{enumitem}
\usepackage{titlesec}
\usepackage{fancyhdr}
\usepackage{xcolor}

\definecolor{castblue}{RGB}{0,48,135}

\pagestyle{fancy}
\fancyhf{}
\renewcommand{\headrulewidth}{0.4pt}
\fancyhead[L]{\small College Area Swim Team}
\fancyhead[R]{\small STP: swim-times}
\fancyfoot[C]{\thepage}

\titleformat{\section}{\large\bfseries\color{castblue}}{\thesection}{0.5em}{}[\titlerule]
\titleformat{\subsection}{\normalsize\bfseries}{\thesubsection}{0.5em}{}

\newcommand{\code}[1]{\texttt{#1}}
\newcommand{\req}[1]{\texttt{#1}}

\begin{document}

% -----------------------------------------------------------------
% Title block
% -----------------------------------------------------------------
\begin{titlepage}
  \centering
  \vspace*{2cm}
  {\Huge\bfseries\color{castblue} Software Test Plan\\[0.5em]}
  {\Large swim-times: Automated Swimmer-Times Retrieval System}\\[2em]
  {\large College Area Swim Team (CAST)}\\[0.5em]
  {\normalsize Prepared for: CAST Coaching Staff and Maintainers}\\[0.5em]
  {\normalsize Prepared by: CAST Technology Working Group}\\[2em]
  \begin{tabular}{ll}
    Document identifier: & CAST-STP-001                                       \\
    Revision:            & 1.0                                                \\
    Date:                & April 22, 2026                                     \\
    Status:              & Released                                           \\
    Governing standard:  & ISO/IEC/IEEE 29119-3:2021 (Test Documentation)     \\
  \end{tabular}
  \vfill
  {\small\textit{This document is prepared in accordance with the principles
  of rational documentation articulated in D.\,L.\,Parnas and P.\,C.\,Clements,
  ``A Rational Design Process: How and Why to Fake It,''
  \textit{IEEE Transactions on Software Engineering}, vol.\,SE-12, no.\,2,
  pp.\,251--257, February 1986.}}
\end{titlepage}

\tableofcontents
\clearpage

% =================================================================
\section{Introduction}
% =================================================================

\subsection{Document identifier}

CAST-STP-001 Rev.\,1.0, \textit{Software Test Plan: swim-times}.

\subsection{Purpose}

This Software Test Plan (STP) specifies how the requirements
defined in CAST-SRS-001 shall be verified by execution.  It
identifies the test items, the test sub-processes, the test
environment, the pass/fail criteria, the deliverables, the
schedule, and the staffing.  The plan is the binding agreement
between development and acceptance regarding what constitutes a
successful test run prior to release.

\subsection{Scope}

This STP covers verification of the swim-times executable that
is tangled out of \code{swim2/swim-times.w}.  It does not cover
verification of upstream USA~Swimming data, the Sisense back-end,
or the unrelated Pandoc-based PDF report at \code{swim/swim-times.md}.

\subsection{References}

\begin{enumerate}[label={[\arabic*]}]
  \item ISO/IEC/IEEE 29119-3:2021, \textit{Software and systems
        engineering---Software testing---Part 3: Test documentation}.
  \item ISO/IEC/IEEE 29148-2018, \textit{Systems and software engineering---
        Life cycle processes---Requirements engineering}.
  \item D.\,L.\,Parnas and P.\,C.\,Clements, ``A Rational Design Process:
        How and Why to Fake It,'' \textit{IEEE Trans.\ Software Eng.},
        vol.\,SE-12, no.\,2, pp.\,251--257, Feb.\ 1986.
  \item CAST-CONOPS-001 Rev.\,1.0, \textit{Concept of Operations:
        swim-times}, April 21, 2026.
  \item CAST-SRS-001 Rev.\,1.0, \textit{Software Requirements
        Specification: swim-times}, April 22, 2026.
  \item CAST-SDD-001 Rev.\,1.0, \textit{Software Design Document:
        swim-times}, April 22, 2026.
  \item \code{swim2/test-swim-times.sh} --- automated test script
        that implements the test cases enumerated in this plan.
\end{enumerate}

% =================================================================
\section{Test Context}
% =================================================================

\subsection{Project background}

swim-times is a single-binary command-line utility that retrieves
swim times from the USA~Swimming data hub via the Sisense JAQL API.
The full operational context is described in CAST-CONOPS-001.

\subsection{Test items}

\begin{itemize}
  \item The CWEB literate source \code{swim2/swim-times.w} (the
        canonical source of record).
  \item The tangled C source \code{swim2/swim-times.c} (a build
        artefact).
  \item The compiled executable \code{swim2/swim-times}.
  \item The cwebmac-typeset documentation PDF
        \code{swim2/swim-times.pdf}.
\end{itemize}

\subsection{Test scope}

\paragraph{In scope.}
All requirements tagged \req{REQ-F-*}, \req{REQ-I-*}, \req{REQ-P-*},
\req{REQ-D-*}, and \req{REQ-A-*} in CAST-SRS-001.

\paragraph{Out of scope.}
\begin{itemize}
  \item The accuracy of data returned by USA~Swimming (treated as
        an authoritative external system).
  \item The cwebmac LaTeX macros and their typesetting choices.
  \item Performance under network outage or DNS failure (these
        are operator-environment conditions).
\end{itemize}

\subsection{Assumptions and constraints}

\begin{itemize}
  \item The test host has outbound HTTPS access to
        \texttt{usaswimming.sisense.com}.
  \item The Sisense bearer token embedded in the source is valid
        on the day the test is executed.
  \item The four tracked swimmers continue to exist in the
        USA~Swimming Persons table.
  \item POSIX utilities (\code{awk}, \code{grep}, \code{diff},
        \code{curl}, \code{otool} or \code{ldd}) are available.
\end{itemize}

\subsection{Stakeholders}

\begin{center}
\begin{tabular}{ll}
  \toprule
  \textbf{Role}            & \textbf{Concern}                              \\
  \midrule
  Maintainer (developer)   & Timely detection of regressions               \\
  Coach (operator)         & Behavioural correctness of CLI                \\
  Auditor                  & Traceability from SRS to evidence             \\
  \bottomrule
\end{tabular}
\end{center}

% =================================================================
\section{Communication}
% =================================================================

Test results are written to standard output by
\code{test-swim-times.sh} as PASS/FAIL/SKIP lines, followed by a
summary count.  A non-zero exit status from the script indicates
test failure and shall be treated as a release blocker by the
maintainer.  No external defect-tracking system is required at
present; failures are recorded as GitHub issues against the
project repository.

% =================================================================
\section{Risk Register}
% =================================================================

\begin{center}
\begin{longtable}{p{4.0cm}p{1.5cm}p{1.5cm}p{6cm}}
  \toprule
  \textbf{Risk} & \textbf{Likeli-} & \textbf{Impact} & \textbf{Mitigation} \\
                & \textbf{hood}     &                  & \\
  \midrule\endhead
  Sisense token rotation invalidates auth
            & Med  & High & Re-embed token, recompile, retest;
                            trigger STP re-run. \\
  Sisense response schema change
            & Low  & High & Test script's \req{REQ-F-32} probe
                            detects parse failure on first run. \\
  Network outage during test
            & Med  & Low  & Online tests \texttt{SKIP} cleanly;
                            offline tests still run. \\
  USA Swimming renames or drops a tracked swimmer
            & Low  & Med  & \req{REQ-F-60} test catches
                            person-lookup failure. \\
  Test host lacks \code{libcurl}
            & Low  & High & Build phase fails fast; STP cannot
                            continue. \\
  \bottomrule
\end{longtable}
\end{center}

% =================================================================
\section{Test Strategy}
% =================================================================

\subsection{Test sub-processes}

Testing is decomposed into four sub-processes, executed in order
by \code{test-swim-times.sh}:

\begin{description}[style=nextline]
  \item[Build sub-process.]
        \code{make tangle compile} produces a fresh executable.
        Failure here aborts the run.
  \item[Static sub-process.]
        Source-code and binary inspection verifies design
        constraints (\req{REQ-D-*}) and self-documentation
        (\req{REQ-A-04}).
  \item[Offline behaviour sub-process.]
        CLI exit-status and usage-text checks (\req{REQ-F-03},
        \req{REQ-F-04}) require no network.
  \item[Online behaviour sub-process.]
        Functional, performance, and reliability checks issue real
        JAQL requests against the Sisense back-end.  Skipped
        cleanly when network is unreachable or the operator sets
        \code{SWIM\_TIMES\_SKIP\_NETWORK=1}.
\end{description}

\subsection{Test design techniques}

\begin{description}[style=nextline]
  \item[Specification-based.]
        Each functional requirement in the SRS yields one or more
        test cases driven by inspection of program output.
  \item[Boundary value.]
        \code{-o} with all four swimmers and \code{-e} with all
        thirty-one default events exercise the upper bound of
        each input domain.
  \item[Equivalence partition.]
        Single-swimmer and single-event invocations represent the
        ``one of many'' equivalence class.
  \item[Static analysis.]
        Compiler warnings (\code{-Wall}), CWEB chunk line-count
        audit, and binary dependency inspection enforce the
        \req{REQ-D-*} constraints.
\end{description}

\subsection{Test environment}

\begin{itemize}
  \item Host: any POSIX.1-2008 system with C99 toolchain,
        \code{libcurl} 7.0+, \code{cweb}, \code{pdftex}, and
        \code{bash} 4+.
  \item Reference development host: macOS 14 with Homebrew
        \code{curl} and TeX~Live~2025.
  \item Network: outbound HTTPS to
        \texttt{usaswimming.sisense.com}.
  \item No databases, message brokers, or auxiliary services are
        required.
\end{itemize}

\subsection{Test data}

The four production swimmer records and the thirty-one production
event codes serve as test data.  No synthetic data is required
because the system under test is read-only against a live
read-only API.

\subsection{Item pass/fail criteria}

\begin{itemize}
  \item A test passes if its \code{PASS} line is printed.
  \item A test fails if its \code{FAIL} line is printed.
  \item A test is skipped (\code{SKIP}) when its preconditions
        cannot be met (e.g., no network).  Skipped tests do not
        block release but are tracked.
  \item The release gate is: zero \code{FAIL} and the count of
        \code{SKIP} not exceeding the online-test count when run
        in offline mode.
\end{itemize}

\subsection{Suspension criteria and resumption requirements}

Testing is suspended if:

\begin{itemize}
  \item The build sub-process fails (cannot proceed without a
        binary).
  \item More than three consecutive online tests fail with
        identical HTTP error codes (suggests a Sisense-side
        outage rather than an SUT defect).
\end{itemize}

Testing resumes once the underlying condition is corrected; the
full STP shall be re-run from scratch, not partially.

\subsection{Test deliverables}

\begin{itemize}
  \item This Test Plan (CAST-STP-001).
  \item The automated test script \code{swim2/test-swim-times.sh}.
  \item Per-run console output captured by the operator
        (recommended: \code{./test-swim-times.sh | tee run.log}).
  \item Updated risk register entries when new failure modes
        are observed.
\end{itemize}

% =================================================================
\section{Test Schedule}
% =================================================================

\begin{center}
\begin{tabular}{ll}
  \toprule
  \textbf{Trigger}                              & \textbf{Action}                  \\
  \midrule
  Any change to \code{swim-times.w}             & Run STP before commit            \\
  Sisense bearer token rotated                  & Run STP after recompile          \\
  USA Swimming schema change announced          & Run STP, update tests if needed  \\
  Quarterly maintenance window                  & Run STP unconditionally          \\
  Pre-release of a tagged version               & Run STP; archive run log         \\
  \bottomrule
\end{tabular}
\end{center}

% =================================================================
\section{Staffing and Roles}
% =================================================================

\begin{center}
\begin{tabular}{ll}
  \toprule
  \textbf{Role}        & \textbf{Responsibility}                                \\
  \midrule
  Maintainer           & Run STP; resolve any \code{FAIL}; update tests when
                         requirements change                                    \\
  Coach (acceptance)   & Verify behavioural tests still match operational need \\
  Reviewer             & Confirm SRS--STP traceability matrix is current       \\
  \bottomrule
\end{tabular}
\end{center}

No external test team or specialised training is required.  The
maintainer who edits \code{swim-times.w} is the same individual
who runs the STP.

% =================================================================
\section{Test Cases and Traceability}
% =================================================================

The table below maps every requirement in CAST-SRS-001 to the
test case in \code{test-swim-times.sh} that verifies it.  Each
row also gives the test method (D = Demonstration, I = Inspection,
T = Test, A = Analysis).

\begin{center}
\begin{longtable}{p{2.6cm}p{1.0cm}p{8.4cm}}
  \toprule
  \textbf{SRS req.} & \textbf{Method} & \textbf{Test case in script} \\
  \midrule\endhead
  \req{REQ-F-01}    & D & Implicit; exercised by every \code{-o} invocation in the suite \\
  \req{REQ-F-02}    & D & Implicit; exercised by every \code{-e} invocation in the suite \\
  \req{REQ-F-03}    & D & Empty-invocation test: \code{rc=1} and \code{Usage:} on stderr \\
  \req{REQ-F-04}    & D & Unknown-flag test: \code{rc=2} and \code{Usage:} on stderr  \\
  \req{REQ-F-10}    & D & \code{-o stella} restricts output to one swimmer            \\
  \req{REQ-F-11}    & D & \code{-e ...} with no swimmer keyword yields four \code{Swimmer:} lines \\
  \req{REQ-F-12}    & D & \code{stella} keyword resolves to a row containing ``Evans'' \\
  \req{REQ-F-20}    & D & Single-event \code{-e} yields exactly one event section    \\
  \req{REQ-F-21}    & D & No \code{-e} yields 31 event sections                      \\
  \req{REQ-F-22}    & D & Event-list count honoured (subset of \req{REQ-F-20})        \\
  \req{REQ-F-30}    & D & Successful baseline run produces a \code{Swimmer:} line     \\
  \req{REQ-F-31}    & D & Successful baseline run produces an event section          \\
  \req{REQ-F-32}    & D & Five output cells render as a complete table row           \\
  \req{REQ-F-33}    & A & Row-cap parameter \code{count:100} verified by source inspection \\
  \req{REQ-F-40}    & T & Insertion-sort fastest-first verified on baseline output   \\
  \req{REQ-F-41}    & D & \code{fastest} mode prints exactly one data row            \\
  \req{REQ-F-42}    & T & All dates match \code{YYYY-MM-DD} regex                    \\
  \req{REQ-F-50}    & D & Heading + column header + separator detected               \\
  \req{REQ-F-51}    & D & CSV header line + at least one data line                   \\
  \req{REQ-F-52}    & D & Empty-event probe yields \code{(no times found)} (else SKIP) \\
  \req{REQ-F-60}    & I & Lookup-failure path inspected in \code{lookup\_person\_key} \\
  \req{REQ-F-61}    & I & Per-event failure-then-continue inspected in \code{fetch\_times} \\
  \req{REQ-I-01}    & I & No \code{stdin} read in source                              \\
  \req{REQ-I-10}    & A & \code{otool}/\code{ldd} confirms libcurl link               \\
  \req{REQ-I-11}    & I & Header strings \code{Content-Type} and \code{Authorization} present in source \\
  \req{REQ-I-20}    & A & \code{https://usaswimming.sisense.com} substring present, no \code{http://} \\
  \req{REQ-P-01}    & D & Sample timing extrapolated to 128 requests $<$ 300 s        \\
  \req{REQ-P-02}    & D & Single query measured $<$ 10 s                              \\
  \req{REQ-D-01}    & A & Compile under \code{-O2 -Wall} produces no warnings         \\
  \req{REQ-D-02}    & I & \code{swim-times.c} regenerable from \code{swim-times.w}    \\
  \req{REQ-D-03}    & A & CWEB chunk line-count audit (awk pass)                     \\
  \req{REQ-D-04}    & A & Binary dependency list contains only \{libc, libcurl, TLS\} \\
  \req{REQ-A-01}    & T & Two consecutive runs produce byte-identical output         \\
  \req{REQ-A-02}    & I & Token and URL declared as named constants                   \\
  \req{REQ-A-03}    & D & Build on a second POSIX host (manual; quarterly)           \\
  \req{REQ-A-04}    & D & \code{cweave} + \code{pdftex} produces \code{swim-times.pdf} \\
  \bottomrule
\end{longtable}
\end{center}

\noindent
Tests marked \textbf{I} (Inspection) are recorded in this plan and
in code review checklists rather than being executed by the
automated script.  All others are executed by
\code{test-swim-times.sh}.

% =================================================================
\section{Approvals}
% =================================================================

\begin{center}
\begin{tabular}{lll}
  \toprule
  \textbf{Role}                & \textbf{Name}            & \textbf{Date}  \\
  \midrule
  Maintainer (author)          & \rule{4cm}{0.4pt}        & \rule{2cm}{0.4pt} \\
  Coach (acceptance)           & \rule{4cm}{0.4pt}        & \rule{2cm}{0.4pt} \\
  Reviewer                     & \rule{4cm}{0.4pt}        & \rule{2cm}{0.4pt} \\
  \bottomrule
\end{tabular}
\end{center}

\vfill
\begin{center}
  \small\textit{End of document CAST-STP-001 Rev.\,1.0}
\end{center}

\end{document}
